\section{Correctness of the Parallel Engine}\label{sec:correctness}

\Cref{thm:samestates} established that the fast path of \cref{sec:filter} does
not change what is explored.  It remains to argue that the parallel engine
itself agrees with \cref{alg:seq}.  We sketch the three obligations; none is
surprising, and we state them mainly to make the trust boundary explicit.

\paragraph{Data-race freedom.}
During expansion, threads read the frontier array and the network, both
immutable for the duration of the level.  The only mutable shared structures are
the passed-list stripes and the zone-intern table.  Every antichain read or
write happens under the owning stripe's mutex, and a discrete key maps to exactly
one stripe, so all operations on a given antichain are totally ordered.  The
intern table serialises its own structural mutations and hands out stable record
addresses (records live in bump-allocated slabs and are never moved), so a reader
holding an index is unaffected by concurrent interning.  Reclamation runs between
levels, when no worker is active.

\paragraph{Verdict equivalence.}
\begin{theorem}[Correctness]\label{thm:correct}
\Cref{alg:par} returns \textbf{true} iff the target $\varphi$ is reachable.
\end{theorem}
\begin{proof}[Sketch]
By induction on the level $k$, with the invariant: after level $k$, every
symbolic state reachable within $k$ steps is either stored in $\Passed$ or
subsumed by a stored one.  The base case is the initial state.  For the step,
every $(q,Z) \in \Frontier_{k-1}$ is expanded, and $\succfn$ is pure, so the
multiset of generated successors is exactly that of \cref{alg:seq} at the same
depth, independently of thread interleaving; each is offered to $\mathrm{add}$,
which either stores it or drops it because a subsuming zone is already present.
Soundness: a \textbf{true} verdict is produced only by $\varphi$ holding at a
genuinely generated state.  Completeness: if a witness $(q^\star,Z^\star)$ is
reachable it is generated at its BFS depth; either it is stored and the target
test fires, or some $Z_p \succeq_{aLU} Z^\star$ was stored earlier, and since
$\preceq_{aLU}$ is sound for reachability~\cite{Herbreteau12} that zone also
witnesses $\varphi$. \qed
\end{proof}

\begin{theorem}[Termination]\label{thm:term}
On diagonal-free networks (or after diagonal splitting), \cref{alg:par}
terminates.
\end{theorem}
\begin{proof}[Sketch]
LU$^+$ extrapolation yields finitely many distinct zones per discrete
state~\cite{Behrmann06}, so $\Passed$ is bounded.  A zone enters a frontier only
when $\mathrm{add}$ stores it, which happens at most once per stored zone.  The
frontiers therefore exhaust and $\Frontier = \emptyset$. \qed
\end{proof}

\paragraph{Non-determinism of the stored set.}
One caveat deserves stating plainly.  Fold~(c) of \cref{alg:add} merges two zones
when their union is exactly convex, and this operation is order-sensitive: which
zones end up merged can depend on the order in which candidates arrive at a
discrete state, and that order depends on thread interleaving.  Consequently the
\emph{exact set of stored zones} is not guaranteed to be bit-identical across
runs, even though \cref{thm:correct} guarantees the verdict is.  In our
experiments the stored-state counts were in fact stable run to run and matched
UPPAAL exactly, but the guarantee we claim is verdict equivalence, not replay
determinism.  This is why the fast path of \cref{sec:filter} is deliberately
placed \emph{before} fold~(c) and phrased as a decision procedure for the
existing predicate: it cannot perturb merge order, and \cref{thm:samestates} is
correspondingly an equality, not an approximation.
